\documentclass[11pt]{article}
\usepackage[T1]{fontenc}
\usepackage{textcomp}
\usepackage[margin=0.75in]{geometry}
\usepackage{amsmath,amssymb,amsthm}
\usepackage{array,booktabs}
\usepackage{hyperref}
\setlength{\parindent}{0pt}
\newtheorem{theorem}{Theorem}
\theoremstyle{definition}
\newtheorem{definition}{Definition}
\title{Flow Proof Artifact: prop-23-no-two-similar-segments-same-side}
\author{Generated by \texttt{flow doc proof}}
\date{Flow Proof Book}
\begin{document}
\maketitle
On the same straight line and on the same side of it there cannot be constructed two similar unequal segments of circles.\\[0.5em]
\textbf{Source.} euclid — Elements, Book III, Proposition 23\\[0.5em]
\bigskip
\noindent{\large \textbf{Derived fact 1.} \textit{On the same straight line and on the same side of it there cannot be constructed two similar unequal segments of circles} \hfill the Euclidean plane \cdot Euclid Book III \cdot Proposition 23: two unequal similar segments cannot stand on the same side of a line}\par
\smallskip\noindent\textit{Coordinate: the Euclidean plane \cdot Euclid Book III \cdot Proposition 23: two unequal similar segments cannot stand on the same side of a line}\par
\smallskip\noindent\textit{Source: euclid — Elements, Book III, Proposition 23}\par
\smallskip\noindent\textit{Built on: proposition 21: angles in the same segment are equal, for Euclid Book III on the Euclidean plane}\par
\medskip
\noindent\textbf{Goal.} On the same straight line and on the same side of it there cannot be constructed two similar unequal segments of circles\par
\noindent$\displaystyle \text{unique segment on line side}$\par
\medskip
\noindent
\begin{tabular}{@{} >{\raggedright\arraybackslash}p{0.06\textwidth} >{\raggedright\arraybackslash}p{0.47\textwidth} >{\raggedleft\arraybackslash}p{0.41\textwidth} @{}}
\textbf{\#} & \textbf{Proof} & \textbf{Mathematics} \\
\midrule
\textbf{1.} & We prove that proposition 23: two unequal similar segments cannot stand on the same side of a line for Euclid Book III the Euclidean plane. & \\[0.45em]
\textbf{2.} & Let AB = straight\_line. & \\[0.45em]
\textbf{3.} & Let seg1 = segment\_on\_AB. & \\[0.45em]
\textbf{4.} & Let seg2 = second\_segment. & \\[0.45em]
\textbf{5.} & From step 2, step 3, and step 4, we can deduce that seg1 equals seg2. & $\displaystyle seg1 = seg2$ \\[0.45em]
\textbf{6.} & We invoke the derived fact governing Euclid Book III the Euclidean plane: proposition 21: angles in the same segment are equal, for Euclid Book III on the Euclidean plane. & \\[0.45em]
\textbf{7.} & From step 6, this implies unique segment on line side. Hence proven. & $\displaystyle \text{unique segment on line side}$ \\[0.45em]
\bottomrule
\end{tabular}
\label{thm:Geometry-Euclid Book III-proposition_23_two_unequal_similar_segments_cannot_stand_on_the_same_side_of_a_line}
\par\medskip\hrule
\end{document}
